\chapter{Where the hands point}
\label{ch:sign-anaphor}

\epigraph{\textit{Between the motion\\And the act\\Falls the Shadow}}{~-- T. S. Eliot, \textit{The Hollow Men} (1925)}

% Chapter 13: Sign language anaphor as HPC test case
% Target: ~6,000 words (7 sections)
% Tests modality independence of HPC architecture
% Core move: anaphoric tracking ≠ spatial agreement

%--- --- --- --- --- --- --- --- --- --- --- --- --- --- --- --- ---
\section{The test}
\label{sec:sign:test}
%--- --- --- --- --- --- --- --- --- --- --- --- --- --- --- --- ---

A signer introduces an absent friend by fingerspelling a name and pointing to a spot in the space to her right. Three sentences later, without repeating the name, she points back to that same location. The hand returns to empty air~-- and the referent\ixsq{referent} is there, even though the friend (inconveniently) is not.

Spoken language does the same work. Pronouns\ixsq{pronoun} track referents across discourse; agreement\ixsq{agreement!in sign languages} marks argument relations on verbs. But much of that machinery is distributed across pronouns, agreement morphology, and constructional constraints. Sign languages\ixsq{sign language} make more of it visible. A locus\ixsq{locus} in signing space is, in some respects, an overt index: you can see referential architecture that spoken language often buries in affix, feature, and clause type.

This chapter asks a direct question. The HPC\ixsq{homeostatic property cluster!sign language} apparatus developed in Part~II~-- clusters, mechanisms, failure modes, projectibility\ixsq{projectibility!in sign languages}~-- was built on spoken-language evidence. Does it survive the switch to a different physical channel? Or does it depend, embarrassingly, on vowels?

The question matters. If the architecture depends on acoustic transmission, it's a claim about spoken-language grammar, not about language as such. If it works when the signal is manual-visual, that's evidence of substrate neutrality\ixsq{modality!independence}~-- the functional organization persists across modalities.

The comparison is fair. Signed and spoken languages share the full structural toolkit: conventional vocabularies, duality of patterning, productivity, syntactic structure, and comparable proposition rates (roughly one every one to two seconds in both modalities; \ixnq{Sandler, Wendy}\ixnq{Lillo-Martin, Diane}\citealt{sandler2006}: 3). What \textcite{sandler2006} call a \enquote{modality independent upper limit}\ixsq{modality!independent upper limit} is supported for the signed and spoken systems they compare (p.~3). Signed languages are independent linguistic systems in a manual-visual modality (\ixnq{Meier, Richard P.}\citealt{meier2002}: 11--12).

The case study is discourse reference\ixsq{discourse reference}~-- the system that introduces, tracks, and reactivates referents across stretches of connected discourse. In spoken language, this system is maintained by pronouns, agreement, word order, and prosody. In sign languages, the same functional job runs through spatial loci, pointing, directional verbs, and body shift. The mechanisms differ, but the architecture is the same.

The chapter parallels the analytical moves of earlier case studies: countability\ixsq{countability} (Chapter~\ref{ch:countability}), where individuation and count syntax turned out to be separable clusters, and definiteness\ixsq{definiteness} (Chapter~\ref{ch:definiteness-and-deitality}), where the semantic notion and the morphosyntactic article system came apart. Here, the same kind of separation appears in a new modality.


%--- --- --- --- --- --- --- --- --- --- --- --- --- --- --- --- ---
\section{How spatial reference works}
\label{sec:sign:how}
%--- --- --- --- --- --- --- --- --- --- --- --- --- --- --- --- ---

Three components make up the spatial reference system in sign languages\ixsq{sign language!spatial reference}. All three run through the same medium~-- the space in front of the signer~-- but they do different work.

First, \term{loci}\ixsq{locus}. Locations in signing space are associated with discourse referents\ixsq{referent}. First and second person are anchored to signer and addressee. Third-person referents get locations established by pointing, naming, or agreement verbs\ixsq{agreement!verbs}. There's no clear upper bound on the number of simultaneously active loci beyond working memory (the grammar is generous, and working memory is less so); \ixnq{Schlenker, Philippe}\textcite{schlenker2017} reports an accepted ASL\ixgq{ASL} example with five third-person loci plus speaker and addressee~-- seven-way tracking (p.~4). In practice, discourse relies on a small active set, typically requiring only two to three recursive spatial splits (\ixnq{Steinbach, Markus}\ixnq{Onea, Edgar}\citealt{steinbach2016}: 2).

Second, \term{pointing}. Sign-language pronominals are indexical points directed at established loci, with both grammatical and gestural structure. They support fine-grained reference tracking through discourse-linked loci. But they aren't purely linguistic. \ixnq{Cormier, Kearsy}\ixnq{Schembri, Adam}\ixnq{Woll, Bencie}\textcite{cormier2013} found that in a BSL\ixgq{BSL} corpus, 789 out of 25,000 pointing tokens were ambiguous in reference and/or function (p.~235). The system tolerates noise. Pronominal pointing \enquote{cannot be characterised exclusively either as pronouns, or as pointing gestures} (\citealt{cormier2013}: 244). Like most things worth studying, it refuses to pick a side.

Third, \term{agreement verbs}\ixsq{agreement!verbs}. A subset of verbs modulate their path or facing to link to argument loci. In ASL\ixgq{ASL}, directional marking is typically stronger toward objects and more variable toward subjects (\citealt{sandler2006}: 30)~-- the verb has its priorities. Verb classes~-- agreeing, plain, spatial~-- are cross-linguistically robust across sign languages, with shared phonetic constraints on realization documented across ASL, DGS\ixgq{DGS}, Auslan\ixgq{Auslan}, and JSL\ixgq{JSL} (\citealt{sandler2006}: 31--38).

A fourth mechanism works alongside these three. \term{Role shift}\ixsq{role shift} (also called constructed action or referential shift) marks perspective change through body position, eye gaze, and facial expression. When a signer \enquote{becomes} a character in a narrative, the entire referential frame shifts: first-person forms now refer to the character, not the signer. Narratives get to borrow the signer for a while. \ixnq{Garcia, Brigitte}\ixnq{Sallandre, Marie-Anne}\textcite{garcia2020} report that these transfer structures account for up to 70\% of reference in sign-language narratives, and dominate in maintenance contexts~-- 76--95\% of cases where a referent is being continued rather than introduced (pp.~6, 8, citing \citealt{pizzuto2011}).

These components have different properties. Pointing is deictic and iconic: the direction of the hand encodes the location of the referent. Agreement verbs are morphological: their path is constrained by verb class and argument structure. Role shift is performative: the signer's body becomes the referent. But they all serve the same functional end~-- keeping track of who's doing what to whom across discourse.

One more property connects the system across modalities. \textcite{steinbach2016} observe that \enquote{the more prominent a discourse referent is, the less explicit the pointing to its R-locus can be} (p.~10), where \enquote{R} labels the discourse referent. This mirrors the accessibility hierarchy familiar from spoken-language reference tracking (\eg Ariel's accessibility theory): full NPs for new referents, pronouns for accessible ones, zero forms for the most prominent. The scale is the same. The forms that realize it differ.


%--- --- --- --- --- --- --- --- --- --- --- --- --- --- --- --- ---
\section{Two clusters, one channel}
\label{sec:sign:two-clusters}
%--- --- --- --- --- --- --- --- --- --- --- --- --- --- --- --- ---

In ASL\ixgq{ASL}, the system looks unified. Referents get loci, pronouns point to loci, and verbs agree with loci. Everything runs through space, so at first pass it can look like a single system.

\textcite{schlenker2017} argues that sign-language loci play \enquote{very much the role of a formal index} (p.~10). They support variable binding, obey crossover constraints, and permit donkey-style anaphoric dependencies (\citealt{zucchi2012}: 2--3). On this view, loci are the discourse-tracking machinery made visible: the same logical operations as in spoken language, realized in space rather than in morphology.

\ixnq{Kuhn, Jeremy}\textcite{kuhn2016} offers a different diagnosis. Same-locus pronouns can receive different interpretations in a single structure. Under the scope of \mention{only one}, two possessive forms pointing to the same locus allow all four readings: bound-bound, bound-free, free-bound, free-free (p.~14--15). If loci were strict variables, same-locus forms would have to co-vary. They don't. \textcite{kuhn2016} concludes that \enquote{the strong loci-as-variables hypothesis has been falsified} (p.~18). Loci pattern instead with morphosyntactic features~-- agreement, underspecification, feature-checking.

\textcite{kuhn2022} pushes this further. Disjunctive referents in LSF\ixgq{LSF} require spatial sums: pointing \enquote{between} two loci retrieves the sum of the referents (p.~15). A single locus can index multiple variables (p.~6). The system needs a featural layer that goes beyond variable assignment: \enquote{loci must be mediated by a featural layer} (p.~1). Even space needs an interface layer.

Both accounts capture real structure, and both agree on key points: loci are grammatically constrained, they participate in discourse reference, and they are not reducible to free co-speech gesture. Schlenker's analysis explains why loci support binding and crossover. Kuhn's explains why they show agreement-like behaviour and tolerate ambiguity. Taken together, these analyses make the same system legible at two resolutions: index-like behaviour for tracking, and feature-mediated behaviour for agreement.

The HPC\ixsq{homeostatic property cluster!sign language} framework makes that complementarity explicit. The literature has long grouped these phenomena under \enquote{spatial reference}; on this account, that umbrella contains two overlapping clusters coupled through the same spatial channel\ixsq{coupling!in sign language} (Table~\ref{tab:sign:two-clusters}).

\begin{table}[htbp]
\caption{Two clusters in sign-language discourse reference.}
\label{tab:sign:two-clusters}
\small
\centering
\begin{tabular}{@{}p{2.8cm}p{5.2cm}p{5.2cm}@{}}
\toprule
 & \textbf{Referent tracking} (semantic) & \textbf{Spatial agreement} (morphosyntactic) \\
\midrule
Function & Introduce, maintain, reactivate discourse referents & Mark participant contrasts (often first vs non-first) and number-related distinctions via directionality; trigger locus features \\
\addlinespace
Maintained by & Spatial working memory, communicative pressure, interactive alignment & Grammatical entrenchment, acquisition regularization, verb-class constraints \\
\addlinespace
Vehicle & Pointing, loci, role shift, gaze & Agreement verbs, person marking, directional morphology \\
\addlinespace
Schlenker's account & Central & Peripheral \\
\addlinespace
Kuhn's account & Peripheral & Central \\
\bottomrule
\end{tabular}
\end{table}

\noindent On this reading, Schlenker foregrounds tracking properties, while Kuhn foregrounds agreement properties. The proposal is not that one account misses the data, but that each emphasizes a different dimension of a coupled system. The locus\ixsq{locus} is the shared spatial vehicle, which is why the two dimensions can look like one system. In mature sign languages like ASL\ixgq{ASL}, the clusters are tightly coupled: loci serve both tracking and agreement simultaneously. But they don't have to be.

This is the same move as Chapter~\ref{ch:definiteness-and-deitality}. There, \term{definiteness}\ixsq{definiteness} (the semantic notion) and \term{deitality}\ixsq{deitality} (the morphosyntactic article-licensing system) were shown to be separable clusters~-- tightly coupled in English, decoupled in languages without articles. Here, \term{anaphoric tracking}\ixsq{anaphoric tracking} (the semantic cluster) and \term{spatial agreement}\ixsq{spatial agreement} (the morphosyntactic cluster) are separable~-- tightly coupled in ASL, decoupled in village sign languages that track referents without agreement morphology.

Evidence for decoupling comes from emergence trajectories in village sign languages and homesign.


%--- --- --- --- --- --- --- --- --- --- --- --- --- --- --- --- ---
\section{What maintains them}
\label{sec:sign:mechanisms}
%--- --- --- --- --- --- --- --- --- --- --- --- --- --- --- --- ---

Four mechanisms\ixsq{mechanism!in sign languages} maintain the referent-tracking cluster. Each contributes to homeostasis: constraining drift, repairing local breakdowns, or re-stabilizing shared conventions over time.

First, spatial working memory\ixsq{mechanism!spatial working memory}. It constrains the number of simultaneously active loci. \textcite{steinbach2016} note that discourse typically supports two to three recursive spatial splits (p.~2). \ixnq{Perniss, Pamela}\ixnq{Ozyurek, Asli@Özyürek, Aslı}\textcite{perniss2015} show a domain-specific contrast: both modalities mark discourse status, but only DGS\ixgq{DGS} predicates showed significantly more spatial localization in reintroduction than maintenance contexts (reintroduction = .70, maintenance = .49), while co-speech gesture predicates showed no such effect (pp.~15--16). The visual modality \enquote{must carry the full load both of identifying referents and marking referential context} (\citealt{perniss2015}: 19).

Second, acquisition\ixsq{mechanism!acquisition}. \ixnq{Morgan, Gary}\textcite{morgan2002} tracked BSL\ixgq{BSL} children's development of complex verb constructions. The youngest group (ages 3--5) produced the full construction in 0\% of cases, defaulting to partial structures~-- the B-part alone~-- in 90\%. By age 9--12, full production reached 70\%. But non-manual marking deficits persisted even in the oldest group (pp.~16--19). \enquote{The developmental problem lies not in understanding the event but in getting the complete conceptual-linguistic mapping} (p.~18). The tracking cluster develops before the morphosyntactic expression catches up.

Third, interactive alignment\ixsq{mechanism!alignment}. Interlocutors converge on shared spatial models through use. \ixnq{Senghas, Ann}\textcite{senghas2001} found that the second cohort of Nicaraguan Sign Language\ixgq{NSL} signers showed higher shared-reference spatial modulation than the first (0.88 vs 0.63, p.~5)~-- not just more spatial marking, but more consistent marking for the same referents. The system becomes more specific and constrained across cohorts, without anyone convening a standards committee.

Fourth, role shift as a co-primary maintenance pathway\ixsq{mechanism!role shift}. When loci decay (too many referents, too much discourse distance), signers shift to constructed action: body lean, eye gaze, facial expression. \textcite{garcia2020} report that these transfer structures dominate in maintenance contexts (pp.~6, 8). This gives the tracking cluster a second robust pathway when locus-based tracking weakens (§\ref{sec:7:what-projection-is}).

Three further mechanisms maintain the agreement cluster\ixsq{spatial agreement} specifically.

Verb-class entrenchment\ixsq{mechanism!entrenchment}. Agreeing, plain, and spatial verb classes are cross-linguistically robust. Morphophonological constraints on their realization hold across four unrelated sign languages (\citealt{sandler2006}: 38). The verb classes are maintained as conventional properties of the lexicon.

Acquisition regularization\ixsq{mechanism!acquisition}. NSL\ixgq{NSL} second-generation signers created consistent directional marking that first-generation signers lacked. \ixnq{Senghas, Ann}\textcite{senghas1997} found that all second-generation signers showed consistent directional layout usage, compared with mixed consistency in the first generation (pp.~9--10). Children regularize the system.

Iconic anchoring. \ixnq{Padden, Carol}\ixnq{Meir, Irit}\ixnq{Aronoff, Mark}\textcite{padden2010} observe that \enquote{the body is a central reference point for all types of verbs} (p.~19). This body-as-subject\ixsq{body-as-subject} anchor provides a stable iconic base from which agreement can grammaticalize. It is, in the relevant sense, hard to misplace.

The two clusters are maintained by partially overlapping but distinguishable mechanism sets. Spatial working memory constrains both; acquisition shapes both but on different timescales; interactive alignment converges the tracking cluster while entrenchment stabilizes the agreement cluster. This mechanism architecture is what makes the clusters separable~-- and what makes their coupling variable across languages.

Taken together, that's a homeostatic profile, not a mere inventory. Under pressure, the system degrades in bounded ways rather than collapsing (limited active splits, locus decay, same-locus ambiguity under overload; \citealt{steinbach2016}: 2, 29--31; \citealt{schlenker2017}: 10--12; \citealt{kuhn2016}: 14--15). It also recruits compensatory repairs~-- explicit locus refresh, role shift, and classifier persistence~-- that restore tracking when one pathway weakens (\citealt{sandler2006}: 91--92; \citealt{garcia2020}: 6, 8). Across cohorts, alignment and acquisition push the system back toward shared conventions (NSL 0.63 to 0.88; \citealt{senghas2001}: 5).


%--- --- --- --- --- --- --- --- --- --- --- --- --- --- --- --- ---
\section{Where they fray}
\label{sec:sign:failure}
%--- --- --- --- --- --- --- --- --- --- --- --- --- --- --- --- ---

The failure modes\ixsq{failure mode!in sign languages} mirror what Chapter~\ref{ch:failure-modes} predicts: maintained clusters degrade in characteristic ways when mechanisms weaken. But from the coupling perspective in Section~\ref{sec:8:coupling}, clusters can also bud and tighten as maintenance strands accumulate.

Over discourse distance, \term{locus decay}\ixsq{failure mode!locus decay} sets in: loci lose specificity, neutral pointing becomes more ambiguous, and prominence plus pragmatics do more resolution work (\citealt{steinbach2016}: 29--31). When a signer points to a location that hasn't been linked to a discourse referent, the result is presupposition failure~-- the sign-language equivalent of a pronoun with no antecedent (\citealt{steinbach2016}: 6). Empty air, sadly, does not count as an antecedent. The same pressure appears as \term{referent competition}\ixsq{failure mode!referent competition} when too many discourse entities are active at once. \textcite{schlenker2017} documents the pattern: same-locus parses degrade, and disambiguation requires finer spatial structure~-- explicit sub-locus distinctions that weren't needed with fewer referents (pp.~10--12). Under \mention{only one}, same-locus possessives allow all four readings (\citealt{kuhn2016}: 14--15). The result is gradient indeterminacy rather than categorical collapse, which is what HPC\ixsq{homeostatic property cluster!failure modes} architecture predicts.

Signers respond with \term{refresh strategies}\ixsq{failure mode!refresh strategies}. Re-establishing a locus (explicit re-pointing) is the simplest repair. Role shift is more expensive but more powerful: the body takes over from the hands, anchoring reference in posture rather than pointing. Classifier predicates offer a third route~-- persistent handshape tracks a referent through an entire predicate chain (\citealt{sandler2006}: 91--92). In Kata Kolok\ixgq{Kata Kolok}, context-poor settings trigger increased use of list buoys and anaphoric strategies over the default absolute pointing (\ixnq{De Vos, Connie}\citealt{devos2012}: 259). This is why \term{ambiguity tolerance}\ixsq{failure mode!ambiguity tolerance} is part of normal operation, not an anomaly. The 789 ambiguous tokens in \citeauthor{cormier2013}'s BSL\ixgq{BSL} corpus aren't system failures. They're the operating margin of a system maintained by redundant mechanisms. Not every token needs to be unambiguous if discourse context, spatial consistency, and interactive alignment can compensate.


%--- --- --- --- --- --- --- --- --- --- --- --- --- --- --- --- ---
\section{The coupling gradient}
\label{sec:sign:gradient}
%--- --- --- --- --- --- --- --- --- --- --- --- --- --- --- --- ---

Section~\ref{sec:8:coupling} introduced a coupling continuum\ixsq{coupling!tightness gradient}: hard, opaque, loose, composite (these labels describe coupling, not relationships). The sign-language evidence now shows that coupling tightness varies \textit{within} a single functional domain across languages~-- just as it does across spoken languages. The tracking and agreement clusters can be tightly linked, loosely linked, or effectively independent.

ASL\ixgq{ASL} sits at the tight end. Both clusters are well-developed and closely linked. Productive agreement verbs, rich locus architecture, binding constraints parallel to spoken-language systems (\citealt{schlenker2017}: 4--5, 10). What you know about the tracking system predicts the agreement system, and vice versa.

BSL\ixgq{BSL} is slightly looser. Spatial agreement is less dominant; body shift and constructed action carry more of the referential load. In \textcite[239]{cormier2013}'s corpus, singular pronominals favour a palm-sideways orientation while plural forms pattern with adverbials~-- the morphosyntactic marking system is less tightly coupled to the discourse-tracking function.

ISL\ixgq{ISL} is intermediate on this coupling gradient~-- looser than tight ASL-like profiles, but tighter than minimal-coupling village-sign profiles~-- and actively changing. \textcite{padden2010} compared three age groups. The youngest cohort showed a divergent spatial pattern (Z-axis 25\%, X-axis 42\%, compared with older groups' Z-dominant profiles) and marked agreement in nearly 75\% of transfer-verb forms. Older cohorts were mostly uninflected~-- only two double-agreement tokens each (pp.~14--18). The agreement cluster is forming in real time, driven by younger signers. Languages do not wait politely for us to finish describing them.

Village sign languages\ixsq{village sign language} show what happens when the clusters decouple.

In \term{Kata Kolok}\ixgq{Kata Kolok} (Bali), the tracking cluster is present: pointing accounts for roughly one in six manual signs in a six-hour corpus (\citealt{devos2012}: 44). But the agreement cluster is largely absent. In 839 tokens of selected transitive predicates, only 3 showed spatial modification~-- far below the threshold for productive agreement (\citealt{devos2012}: 152--153). Kata Kolok uses an absolute spatial frame rather than body-relative loci. \enquote{Exophoric reference predominates} (\citealt{devos2012}: 259). The tracking cluster remains robust without agreement morphology.

In \term{ABSL}\ixgq{ABSL} (Al-Sayyid Bedouin Sign Language, Israel), word order and prosodic marking emerged early; agreement didn't. Older signers frequently avoided multi-argument clauses: 47\% of clauses used a single argument (\ixnq{Sandler, Wendy}\ixnq{Meir, Irit}\ixnq{Padden, Carol}\ixnq{Aronoff, Mark}\citealt{sandler2005}: 6). \enquote{The transition into language is gradual and not abrupt} (\citealt{sandler2005}: 26). No agreement morphology was reported after multiple generations (p.~28). Distinct subsystems stabilize at different rates~-- syntax and prosody before agreement.

In \term{NSL}\ixgq{NSL} (Nicaraguan Sign Language), emergence has been tracked across cohorts since the 1970s. First-generation signers didn't produce two-animate-arguments-with-one-verb structures; they split events into one-argument clauses with interleaved verbs (\citealt{senghas1997}: 5--6). Second-generation signers introduced consistent directional layout and new argument orders (p.~9). But productive agreement is still developing. \enquote{Not even a highly motivated grammatical system like agreement emerges overnight} (\citealt{padden2010}: 2). Grammaticalization keeps its own clock.

In \term{homesign}\ixsq{homesign} systems, \ixnq{Goldin-Meadow, Susan}\ixnq{Feldman, Heidi}\textcite{goldinmeadow1977} studied six deaf children of hearing parents who received no sign-language input. The children built lexicons with large deictic inventories~-- 41--62\% of each child's signs were deictic (p.~2). They produced systematic semantic ordering in two-sign phrases (pp.~3--4). These data show tracking resources without evidence of productive agreement morphology.

The pattern across these languages supports a tendency rather than a universal sequence. Across the systems reviewed here, tracking resources stabilize early and broadly, including in homesign systems. Productive agreement morphology follows more variable, ecology-sensitive trajectories. The agreement cluster can bud from the tracking system through grammaticalization, just as the article system budded from demonstratives in Chapter~\ref{ch:definiteness-and-deitality}. In both cases, the semantic function existed first. A morphosyntactic mechanism grammaticalized to serve it. Once established, that mechanism developed its own internal constraints. And the two can decouple\ixsq{budding}.

Table~\ref{tab:sign:gradient} summarizes the coupling gradient across sign languages.

\begin{table}[htbp]
\caption{Coupling gradient across sign languages.}
\label{tab:sign:gradient}
\small
\centering
\begin{tabular}{@{}p{1.6cm}@{\hspace{0.35em}}p{4.6cm}p{4.5cm}p{1.3cm}@{}}
\toprule
\textbf{Language} & \textbf{Tracking evidence} & \textbf{Agreement evidence} & \textbf{Coupling} \\
\midrule
ASL & Productive locus reuse and binding constraints (\citealt{schlenker2017}: 4--5, 10) & Productive directional marking and participant contrasts (\citealt{sandler2006}: 30--38) & tight \\
BSL & Robust pointing and locus use; 789/25,000 pointing tokens ambiguous in reference/function (\citealt{cormier2013}: 235) & Palm-orientation marking differentiates singular/plural patterns and aligns plurals with adverbials (\citealt{cormier2013}: 239) & moderate \\
ISL & Discourse-level spatial tracking and referent localization are established (\citealt{padden2010}: 14--18) & Young cohort marks transfer-verb agreement in nearly 75\%; older cohorts mostly uninflected (\citealt{padden2010}: 14--18) & intermediate \\
NSL & Shared-reference spatial modulation rises across cohorts (0.63 to 0.88; \citealt{senghas2001}: 5) & Consistent directional layout appears in later cohorts; productive agreement still developing (\citealt{senghas1997}: 9--10) & loose \\
ABSL & Discourse tracking available via word order, clause sequencing, and prosody (\citealt{sandler2005}: 6, 26) & No agreement morphology reported after multiple generations (\citealt{sandler2005}: 28) & minimal \\
Kata Kolok & Pointing is frequent (roughly 1 in 6 manual signs; \citealt{devos2012}: 44) & Spatial modification is rare (3/839 selected transitive tokens; \citealt{devos2012}: 152--153) & minimal \\
Homesign & Deictic inventories are large (41--62\% of signs; \citealt{goldinmeadow1977}: 2) & No productive agreement morphology is described in these data (\citealt{goldinmeadow1977}: 3--4) & tracking only \\
\bottomrule
\end{tabular}
\end{table}

\noindent Across the systems reviewed here, pointing-based tracking is robust. Locus architecture and agreement morphology vary with ecology, transmission history, and conventionalization, not time depth alone. More maintenance strands generally produce tighter coupling; fewer strands produce looser coupling and greater reliance on context. This is the same prediction the book makes for spoken languages.


%--- --- --- --- --- --- --- --- --- --- --- --- --- --- --- --- ---
\section{The verdict}
\label{sec:sign:verdict}
%--- --- --- --- --- --- --- --- --- --- --- --- --- --- --- --- ---

The verdict turns on three questions.

Is there a property cluster\ixsq{homeostatic property cluster!sign language}? Yes. Referent introduction, maintenance, reactivation, binding, accessibility decay, agreement marking, person distinction, role shift~-- these co-occur reliably in mature sign languages, but not invariably. Village sign languages show subsets. The cluster is real and decomposable.

Are there homeostatic mechanisms\ixsq{mechanism!in sign languages}? Yes. The argument is regulatory, not taxonomic: working-memory limits bound load, repair strategies recover tracking after degradation, and alignment plus acquisition re-stabilize conventions across interaction and cohorts (\citealt{steinbach2016}: 2, 29--31; \citealt{senghas2001}: 5). Verb-class entrenchment and body anchoring further stabilize agreement contrasts where agreement is productive (\citealt{sandler2006}: 38; \citealt{padden2010}: 19). These mechanisms are specific, testable, and modality-sensitive~-- spatial working memory rather than phonological form, body-as-subject rather than morphological paradigm~-- but functionally parallel to spoken-language mechanisms.

Does it project\ixsq{projectibility!in sign languages}? Yes. Knowing that a sign language has productive agreement verbs supports probabilistic predictions about tighter locus reuse, richer participant-marking contrasts, and more agreement-linked null arguments. Knowing it lacks such morphology shifts expectations toward pointing-based tracking, role shift, and discourse prominence for resolution (\citealt{sandler2006}: 15--16; \citealt{steinbach2016}: 6, 10). The cluster earns its inferential keep.

The modality-independence\ixsq{modality!independence} claim holds~-- with a structured caveat. The architecture is substrate-neutral: clusters maintained by mechanisms, projecting for purposes. The mechanisms themselves are modality-sensitive. Spatial working memory replaces phonological entrenchment as a key constraint; the body-as-subject anchor has no spoken-language equivalent. But the \textit{organization}~-- clusters stabilized by causal mechanisms, degrading in characteristic ways when mechanisms weaken, projecting reliably when the cluster is maintained~-- is the same.

What this chapter adds to the argument of Part~III is a modality test rather than another spoken-language case study. Countability, definiteness, lexical categories, and gender (Chapters~\ref{ch:countability}--\ref{ch:proform-gender}) all demonstrated the HPC architecture in spoken language. Sign-language discourse reference demonstrates it in a different physical medium, with different mechanisms maintaining the same functional organization.

The central result can be stated directly. The sign-language literature already identifies multiple functions under the umbrella label \enquote{spatial reference}; this chapter systematizes that evidence as two overlapping clusters: anaphoric tracking\ixsq{anaphoric tracking} and spatial agreement\ixsq{spatial agreement}. The Schlenker/Kuhn debate is more tractable once both are treated as real clusters. Evidence from village sign languages shows that the clusters can decouple. The agreement system budded from the tracking system, just as the article system budded from demonstratives. Coupling tightness varies across sign languages for the same mechanistic reasons it varies across spoken languages: more maintenance strands, tighter coupling.

Where does this leave the coupling spectrum from Section~\ref{sec:8:coupling}? Sign-language discourse reference sits in the loose-to-moderate range. The tracking cluster has transparent mechanisms and predictable failure modes, like gender (Chapter~\ref{ch:proform-gender}). The agreement cluster is more variable~-- tightly coupled in some languages, absent in others. The domain is narrower than grammatical categories but wider than a single construction.

The next chapter (Chapter~\ref{ch:the-category-zipper}) maps this spectrum explicitly: from phonemes, where coupling is hard and failures are local, through words, grammatical categories, and constructions, where coupling loosens and the interpretant does progressively more mediating work. Sign-language discourse reference provides a new data point on that gradient~-- one where the mechanism set differs, but the architecture persists.
